%%%PAGE 150 rot=0 legibility=good summary=φ-x图像(斜率k=-E_x)与E-x图像(面积为U)；静电感应/平衡/屏蔽(近异远同、内部场强为0、导体等势体、内部挖空静电屏蔽)及法拉第笼等示意图；磁感应强度定义式；洛伦兹力分力可做功例题(磁场移动使木块离地)
% 页面上另有一条从底部“V=mg/qB”附近斜向上延伸到电感应示意图下方的手绘长线，未转录；页眉“No./Date.”及背面透印字迹忽略
%%%BLOCK id=150-01 ch=09 sec=E-x与φ-x图像 kind=knowledge alt=none cont=none y=0.0-0.40
% 上图为 φ-x 曲线(图内标有 A、B 点，B 点左侧“k>0，E向左”，右侧“k=-E_x”，两侧各画一个 E 方向箭头)；下图为 E-x 曲线(图内标有 E、x、0、A、U 及弯箭头)
%FIG p150_a 0.17 0.004 0.55 0.19
\notefig[0.3]{figures/p150_a.png}
\noindent $E$-$x$　　　　　$\varphi$-$x$

\noindent 面积为$U$
%FIG p150_b 0.08 0.248 0.24 0.382
\notefig[0.25]{figures/p150_b.png}
\noindent $E_x$：与$x$轴平行（$E$-$x$图）% E_x 后的标点辨认不清(冒号或句点)

\noindent 面积不用带负号
%%%END
%%%BLOCK id=150-02 ch=09 sec=静电平衡 kind=summary alt=none cont=none y=0.17-0.52
% 红笔：导体矩形下方的两条向左红色箭头(感应电荷产生的场)及红字“到抵消时”；(2)下有下划线；(3)的“电场线处处垂直等势面”下有长下划线；(4)下有下划线；原稿(3)分三行书写：“导体为等势体——表面为等 / 势面——电场线处处垂直等 / 势面”
\textbf{静电感应/平衡/屏蔽}

\noindent 过程　　结果　　应用
%FIG p150_c 0.53 0.24 0.70 0.335
\notefig[0.25]{figures/p150_c.png}
\noindent \red{\unclear{到}\,抵消时}

\noindent (1) 近异远同

\noindent (2) \underline{内部无净电荷场强为$0$}

\noindent (3) 导体为等势体——表面为等势面——\underline{电场线处处垂直等势面}

\noindent (4) \underline{内部挖空：静电屏蔽}
%FIG p150_d 0.455 0.36 0.695 0.51
\notefig[0.25]{figures/p150_d.png}
%%%END
%%%BLOCK id=150-03 ch=09 sec=静电屏蔽 kind=diagram alt=none cont=none y=0.40-0.52
% 图中标注“法拉第笼”，引线指向内部方框(方框内画有一个人形)；笼壁左侧标“+”、右侧标“=”；图右侧“⊕”为外部正电荷
%FIG p150_e 0.10 0.405 0.36 0.515
\notefig[0.25]{figures/p150_e.png}
%%%END
%%%BLOCK id=150-04 ch=09 sec=静电平衡 kind=knowledge alt=none cont=none y=0.52-0.63
%FIG p150_f 0.09 0.515 0.255 0.625
\notefig[0.25]{figures/p150_f.png}
净电荷分布在导体外表面
%%%END
%%%BLOCK id=150-05 ch=09 sec=静电平衡 kind=diagram alt=none cont=none y=0.48-0.71
% 两幅静电感应示意图：各画两个导体球(带支架，球上用红笔圈出感应电荷)、接地符号及一个外部⊕电荷(以长线连出)，无文字
%FIG p150_g 0.52 0.475 0.855 0.595
\notefig[0.33]{figures/p150_g.png}
%FIG p150_h 0.52 0.592 0.78 0.71
\notefig[0.26]{figures/p150_h.png}
%%%END
%%%BLOCK id=150-06 ch=11 sec=磁感应强度·定义式 kind=formula alt=none cont=none y=0.72-0.75
\noindent $B=\dfrac{F}{IL\sin\theta}$　　电生磁　一\unclear{速}二\unclear{垂}三右手% “速”“垂”两字笔迹难辨
%%%END
%%%BLOCK id=150-07 ch=11 sec=洛伦兹力分解·分力做功 kind=problem alt=06 cont=none y=0.75-0.90
% 图中红笔标注：F洛1、V1、V合、F洛2 等；黑笔有 qVB、V、mg；图为地面上的小木块，其周围画有“×”(磁场垂直纸面向里)
\noindent 洛伦兹力　分力可以做功，合力不做功

\begin{problem}
%FIG p150_i 0.045 0.776 0.225 0.88
\notefig[0.25]{figures/p150_i.png}
使小木块对地分离　求磁场移动速度
\begin{solution}
$qVB=mg$　使$B$向左以$V=\dfrac{mg}{qB}$运动
\end{solution}
\end{problem}
%%%END
%%%PAGEEND 150
